\documentclass{article}
\usepackage[left=2cm, right=2cm, top=2cm, bottom=2cm]{geometry}
\usepackage{graphicx}
\usepackage{amsmath}
\usepackage{amssymb}
\usepackage{amsthm}
\usepackage{fancyhdr}
\usepackage{verbatim}
\usepackage{listings}
\usepackage{xcolor}
\usepackage{pdfpages}
\usepackage{cancel}
\usepackage{physics}

\lstset{
    language=C++,
    basicstyle=\ttfamily\footnotesize,
    keywordstyle=\color{blue},
    commentstyle=\color{green},
    stringstyle=\color{red},
    numbers=left,
    numberstyle=\tiny\color{gray},
    frame=single,
    breaklines=true,
    captionpos=b,
}


\title{MTH 553: Lecture \# 25}
\author{Cliff Sun}

\newtheorem{theorem}{Theorem}[section]
\newtheorem{lemma}[theorem]{Lemma}
\newtheorem{definition}[theorem]{Definition}
\newtheorem{conjecture}[theorem]{Conjecture}
\newtheorem{proposition}[theorem]{Proposition}
\newtheorem{corollary}[theorem]{Corollary}
\newtheorem{one minute paper}[theorem]{One Minute Paper}

\pagestyle{fancy}
\lhead{\textbf{\thepage}\ \ \nouppercase{\rightmark}}
\chead{MTH 553: Lecture \# 25}
\rhead{Cliff Sun}

\begin{document}

\maketitle

\section*{Lecture Span}
\begin{enumerate}
    \item Finish heat eqiation
    \item Start laplace equation
\end{enumerate}

\begin{theorem}
    Uniqueness, backwards in time for heat equation. Let $\Omega \subset \mathbb{R}^n$ be bounded and smooth, then there is a one solution $v \in C^{2}(\overline{\Omega}_T)$.
\end{theorem}

In other words, 

\begin{align*}
    u_t - \triangle u &= f\\
    u &= g, \quad \text{on the sides and top}
\end{align*}

Note that we are specifiying BC on the sides and on the top. Then we can reverse everything. 

\begin{proof}
    Suppose $u_1,u_2$ are two solutions. Let $w = u_1 - u_2$ so then 
    \begin{align*}
        w_t - \triangle w &= 0\\
        w &= 0, \quad \text{on the sides and top}
    \end{align*}
    Let 
    \begin{equation}
        E(t) \equiv \int_{\Omega}w(x,t)^2dt
    \end{equation}
    Note that $E(T) = 0$. We want to obtain that $E(t) =0 $ for all $t$. Then recall 
    \begin{align*}
        E'(t) &= -2\int_{\Omega}ww_t dx \\
        &= -2\int_{\Omega}w\triangle w dx \\
        &= -2 \int |\nabla w|^2 dx \leq 0
    \end{align*}
    Next, we differentiate once again, 
    \begin{align*}
        E''(t) &= -4\int_{\Omega}\nabla w \cdot \nabla w_t dx \\
        &= 4 \int_{\Omega}\triangle w w_t dx, \quad \text{by divergence theorem (IBP) using $w=0$ on sides, therefore $w_t =0$ on sides.} \\
        &= 4\int_{\Omega}(\triangle w)^2 dx \geq 0
    \end{align*}
    Therefore, our energy is always going down in a "convex" fashion. But this isn't enough to state that $E(t) =0$ for all $t$. 
    Next, we apply Cauchy-Schwartz. 
    \begin{align*}
        \implies E'(t)^2 &= 4 \left( \int_{\Omega}w \triangle w dx \right) \\
        &\leq 4 \int_{\Omega}w^2 dx \cdot \int_{\Omega}(\triangle w)^2 dx \\
        = 4E(t)E''(t)
    \end{align*}
    Case 1, if $E(t) = 0$ for all $t$, then we are done. Case 2, otherwise, there exists some interval where $E(t)$ is positive. In particular, choose some $[t_1,t_2]$ such that $E(t_2)$ is the point of which the $E(t_2)$ becomes zero.
    Then, take $\log(E(t))''$, then obtain 
    \begin{equation}
        \left( \frac{E'(t)}{E(t)} \right)' = \frac{E''(t)}{E(t)} - \frac{E'(t)^2}{E(t)^2} \geq 0
    \end{equation}
    This is because we proved that $E'(t)^2 \leq E(t)E''(t)$. So $\log(E(t))$ is convex. But we know that $\log(E_2(t)) \rightarrow -\infty$ as $t \rightarrow t_2$. Therefore, this cannot be convex. This is because a convex function must come up at some point. 
    This means that case 2 cannot happen. Therefore, we are in case 1 and we are done.
\end{proof}

\section*{Laplace's Equation}

Call $u$ harmonic if 
\begin{equation}
    \triangle u = 0
\end{equation}

This is \underline{Laplace's Equation}. If $u$ is at equilibrium temperature, that is it doesn't depend on time. Then 
\begin{equation}
    u_t = \triangle u \implies \triangle u = 0
\end{equation}

This is a steady state temperature. Poisson's equation is a steady state equilibrium temperature with a driving force. That is 
\begin{equation}
    -\triangle u = f
\end{equation}
Here $f$ is a driving force. Here, if $f > 0$, then it is a heat source. vice versa means heat sink. Elemental harmonic functions
\begin{enumerate}
    \item Linear, $u(x) = a_1x_1 + a_2x_2 + \dots + b$
    \item $n=2$, $u=x_1^2 - x_2^2$. This is a saddle shape, and is generic for graphs of harmonic functions.
    \item $n=2$, $u= 1/\log|x| = 1/2 \cdot \log(x_1^2 + x_2^2)$. This is the \underline{fundamental} solution of the Laplacian. Here, the graph has a saddle shape in the $x_1$ and $x_2$ direction. 
    \item $n=3$, then $u(x) = 1/|x|$ is also harmonic. Here, $|x| = \sqrt{\sum^3 x_i^2}$. 
\end{enumerate}

\end{document}